% !TeX program = lualatex
\documentclass[11pt, a4paper]{report}

% --- UNIVERSAL PREAMBLE BLOCK ---
\usepackage[a4paper, top=2.5cm, bottom=2.5cm, left=2.5cm, right=2.5cm]{geometry}
\usepackage{fontspec}

\usepackage[spanish, bidi=basic, provide=*]{babel}

\babelprovide[import, onchar=ids fonts]{spanish}
\babelprovide[import, onchar=ids fonts]{english}

% Set default/Latin font to Serif in the main (rm) slot for an academic look
\babelfont{rm}{Noto Serif}
% Assign a specific font for Spanish text
\babelfont[spanish]{rm}{Noto Serif}

% Add because main language is not English
\usepackage{enumitem}
\setlist[itemize]{label=-}
% --- END UNIVERSAL PREAMBLE BLOCK ---

\usepackage{amsmath}
\usepackage{booktabs}
\usepackage{titlesec}
\usepackage{setspace}

% Formato de los capítulos para que sean más estéticos
\titleformat{\chapter}[display]
  {\normalfont\bfseries\Huge}{\chaptertitlename\ \thechapter}{20pt}{\Huge}

% Hyperref siempre al final
\usepackage[colorlinks=true, linkcolor=blue, citecolor=black, urlcolor=blue]{hyperref}

\begin{document}

\begin{titlepage}
    \centering
    \vspace*{2cm}
    
    {\Large \textbf{TRABAJO DE FIN DE GRADO}} \\[1.5cm]
    
    {\huge \textbf{Orchestration Trade-offs in Edge Computing: Systemd vs Kubernetes}}\\[0.5cm]
    {\Large Análisis de rendimiento, distribución de cargas e inferencia de modelos de IA en el extremo de la red}\\[2cm]
    
    {\Large \textbf{Autor:}} \\
    {\large [Tu Nombre y Apellidos]}\\[1cm]
    
    {\Large \textbf{Director/es:}} \\
    {\large [Nombre del Director/a]}\\[2cm]
    
    {\large \textbf{Grado en Ingeniería [Informática / Telecomunicaciones]}} \\[0.5cm]
    {\large [Nombre de tu Universidad]} \\[0.5cm]
    {\large [Mes y Año de la convocatoria]}
    
    \vfill
\end{titlepage}

\pagenumbering{roman}
\tableofcontents
\listoffigures
\listoftables
\cleardoublepage

\pagenumbering{arabic}
\setstretch{1.15}

\chapter{Introducción}

Este capítulo debe introducir al lector al problema, por qué es importante y qué vas a resolver.

\section{Contexto y Motivación}
Aquí debes hablar sobre el auge del \textit{Edge Computing} (computación en el extremo) y cómo difiere del \textit{Cloud Computing} tradicional. Explica que los entornos Edge tienen recursos limitados (CPU, RAM, red inestable) y que las herramientas de orquestación diseñadas para la nube (como Kubernetes) pueden ser demasiado pesadas. Introduce el dilema: ¿Necesitamos orquestadores complejos o bastan herramientas nativas del sistema operativo como Systemd + Podman?

\section{Objetivos del Trabajo}
\begin{itemize}
    \item \textbf{Objetivo Principal:} Analizar y comparar los \textit{trade-offs} (rendimiento, resiliencia y complejidad) entre Kubernetes y Systemd como orquestadores para cargas de trabajo Edge.
    \item \textbf{Objetivos Específicos:}
    \begin{itemize}
        \item Desplegar una infraestructura Master-Worker utilizando Ansible, Podman y Systemd (Fase 1).
        \item Evaluar el rendimiento de múltiples protocolos de red (HTTP, gRPC, ZeroMQ, MQTT) y formatos de serialización (JSON, Protobuf, MessagePack) para la transmisión masiva de datos (Dataset MNIST).
        \item Implementar un flujo de trabajo de Machine Learning distribuido (Fase 2) entrenando un modelo PyTorch en el nodo máster e infiriendo en los nodos worker.
        \item Migrar la arquitectura a un entorno de Kubernetes/K3s (Fase 3) y contrastar las métricas de consumo y latencia frente a Systemd.
    \end{itemize}
\end{itemize}

\section{Estructura de la Memoria}
Breve resumen de lo que contiene cada uno de los siguientes capítulos.

\chapter{Estado del Arte y Tecnologías Involucradas}

Define teóricamente las herramientas y conceptos que usas. No es un manual de uso, sino una justificación técnica.

\section{Edge Computing y Orquestación}
Diferencias entre Cloud, Fog y Edge. Desafíos en el Edge (huella de memoria, latencia, ancho de banda). 

\section{Gestión de Contenedores}
\begin{itemize}
    \item \textbf{Podman:} Arquitectura \textit{daemonless} y \textit{rootless}. Por qué es mejor que Docker para el Edge.
    \item \textbf{Systemd Quadlets:} Cómo Systemd integra de forma nativa la gestión del ciclo de vida de los contenedores.
    \item \textbf{Kubernetes (K3s):} Distribuciones ligeras de Kubernetes para IoT/Edge. Arquitectura del plano de control y el CNI (Container Network Interface).
\end{itemize}

\section{Protocolos de Comunicación y Serialización}
Explica la teoría y diferencias de los que has comparado:
\begin{itemize}
    \item \textbf{REST / HTTP(s) y JSON:} Estándar web, basado en texto. Cuellos de botella en el parseo.
    \item \textbf{gRPC y Protocol Buffers:} Llamadas a procedimientos remotos sobre HTTP/2. Serialización binaria estructurada y fuertemente tipada.
    \item \textbf{ZeroMQ:} Sockets TCP crudos sin \textit{broker} (Brokerless P2P). Alta velocidad.
    \item \textbf{MQTT:} Patrón Pub/Sub con \textit{broker} intermediario.
    \item \textbf{MessagePack:} Serialización binaria sin esquema (\textit{schema-less}), flexible y rápida.
\end{itemize}

\section{Machine Learning Distribuido}
Breve introducción a PyTorch, inferencia de modelos, y el dataset MNIST.

\chapter{Diseño del Sistema y Entorno de Pruebas}

\section{Infraestructura Física y Virtual}
Detalla el hardware:
\begin{itemize}
    \item Nodo Master (Padre): Ubuntu 22.04.5 bajo WSL en Windows 11. Rol: Orquestador (Ansible), entrenamiento de IA, y despachador de tareas.
    \item Nodos Workers (Hijos): 2 Máquinas Virtuales Ubuntu 24.04.4 en VMware. Rol: Inferencia, cálculo, retorno de métricas.
\end{itemize}

\section{Aprovisionamiento y Despliegue Agnóstico}
Explica la separación estricta entre \textit{Build} (construcción en el máster) y \textit{Run} (ejecución en el worker). Detalla cómo Ansible transporta los artefactos locales (`.tar`) usando `podman load` (estrategia \textit{side-loading}) para ahorrar ancho de banda y CPU, evitando repositorios externos (Docker Hub).

\section{Arquitectura Lógica de la Solución}
Aquí debes poner el esquema de la arquitectura Master-Worker, mostrando cómo fluyen los datos (Fase 2, Fase 3, Fase 4).

\begin{figure}[htbp]
  \centering
  \framebox{\parbox{0.8\textwidth}{\centering
    \vspace{3cm}
    \textbf{Esquema de Arquitectura Master-Worker} \\
    \small\textit{Aquí se insertará tu diagrama (flow.png) mostrando el nodo Padre, los hijos y el flujo de red.}
    \vspace{3cm}
  }}
  \caption{Topología del sistema y flujos de comunicación.}
  \label{fig:arquitectura}
\end{figure}

\chapter{Fase 1: Transmisión de Datos en Systemd (Conteo de Imágenes)}

En este capítulo desarrollas todo lo que me has pasado sobre el conteo de imágenes y la comparativa de red.

\section{Definición de la Tarea}
El Master descarga el dataset MNIST (60,000 imágenes), lo carga en memoria, lo divide en $N$ particiones ($tamaño\_particion = total\_imagenes / N\_nodos$) y envía los lotes. El Worker cuenta las imágenes recibidas y devuelve el total al Master.

\section{Diseño del Experimento (Comparativa)}
Explica la justificación de cada prueba:
\begin{itemize}
    \item \textbf{Línea Base:} HTTP/1.1 + JSON (Demostrar ineficiencia del texto plano).
    \item \textbf{Estándar Microservicios:} gRPC + Protobuf (HTTP/2 binario).
    \item \textbf{Rendimiento Crudo:} ZeroMQ + Protobuf (Límite físico P2P).
    \item \textbf{IoT Broker:} MQTT + Protobuf (Penalización por doble salto).
    \item \textbf{Flexibilidad:} ZeroMQ + MessagePack (Esquema estático vs dinámico).
\end{itemize}

\section{Resultados y Análisis de Rendimiento}
Inserta aquí las tablas y gráficas de la Fase 1. Comenta el consumo extremo de RAM del JSON frente al paradigma binario (\textit{Zero-Copy} con `np.frombuffer`).

\begin{table}[htbp]
    \centering
    \begin{tabular}{llccc}
        \toprule
        \textbf{Protocolo} & \textbf{Formato} & \textbf{Tiempo RTT (ms)} & \textbf{Throughput (img/s)} & \textbf{RAM Pico (MB)} \\
        \midrule
        HTTP/REST & JSON & $\approx 16500$ & 3,636 & 2600 \\
        gRPC & Protobuf & $\approx 1037$ & 57,850 & 384 \\
        ZeroMQ & Protobuf & $\approx 672$ & 89,285 & 250 \\
        ZeroMQ & MsgPack & $\approx 540$ & 111,087 & 209 \\
        \bottomrule
    \end{tabular}
    \caption{Resumen de resultados de transmisión para 60,000 imágenes.}
    \label{tab:resultados_red}
\end{table}

\section{Conclusiones de la Fase 1}
Explica cómo "La tríada ZeroMQ + MessagePack + Zero-Copy se establece como la arquitectura ganadora" dentro de Systemd.

\chapter{Fase 2: Despliegue e Inferencia de Modelos de IA}

Este capítulo evoluciona el envío tonto de imágenes a una tarea computacional real.

\section{Entrenamiento del Modelo en el Master}
Explica cómo el nodo Master utiliza PyTorch para entrenar un modelo clasificador convolucional (CNN) sencillo sobre el dataset MNIST. Detalla las métricas del entrenamiento (Epochs, Loss, Accuracy final).

\section{Distribución del Modelo y Datos}
Explica cómo se envía el modelo entrenado (ej. serializado con `torch.save` en formato `.pt` o `.pth`) hacia los workers utilizando la arquitectura de red ganadora de la Fase 1 (ZeroMQ). Seguidamente, se envían las particiones de imágenes.

\section{Ejecución de Inferencia en el Edge (Workers)}
Los nodos hijos reciben el modelo, lo cargan en memoria (`torch.load`), y para cada imagen recibida realizan una inferencia (predicción del número escrito a mano). 

\section{Métricas de Procesamiento ($T_{proc}$)}
\begin{itemize}
    \item Tiempo medio de inferencia por lote.
    \item Consumo de CPU de los workers durante la carga computacional (Tensores IA).
    \item Fiabilidad: Exactitud de los números predichos y devueltos al máster.
\end{itemize}

\chapter{Fase 3: Migración a Kubernetes (K3s) y Análisis de Trade-offs}

Aquí culmina el propósito del TFG: la comparación directa de orquestadores.

\section{Migración de la Infraestructura a K3s}
Explica la configuración del cluster Kubernetes. 
\begin{itemize}
    \item Master asume el rol de Control Plane.
    \item Workers asumen roles de Kubelets.
    \item Traducción de los servicios `Systemd Quadlets` a manifiestos de Kubernetes (`Deployments`, `Services`, `Pods`).
\end{itemize}

\section{Ejecución de Cargas de Trabajo (Fase 1 y 2 en K8s)}
Se repite el envío masivo de imágenes y la inferencia de IA, pero ahora a través de la red virtual de Kubernetes (ej. Flannel o Calico CNI) y Kube-proxy.

\section{Comparativa Directa: Systemd vs Kubernetes}
Aquí comparas los resultados de Systemd frente a K8s usando las métricas definidas:
\begin{itemize}
    \item \textbf{Plano de Gestión (Reposo):} Consumo de RAM y CPU de Systemd+Podman inactivo frente al Kubelet+Kube-proxy+Etcd inactivos en el Edge.
    \item \textbf{Plano de Datos (Red):} Cuánta latencia extra añade el CNI de K8s a gRPC y ZeroMQ frente a la comunicación nativa P2P de Podman.
    \item \textbf{Complejidad Operativa:} Facilidad de despliegue usando Ansible estático vs. mantenimiento del estado en K8s.
\end{itemize}

\chapter{Conclusiones y Trabajo Futuro}

\section{Conclusiones Finales}
Resumen general. Responde a la pregunta del TFG: ¿Vale la pena el peso de Kubernetes en el Edge, o es Systemd una alternativa viable? 
Destaca los números clave: 
1) La inviabilidad de REST/JSON. 
2) La eficiencia extrema de ZeroMQ+MsgPack en Podman. 
3) El *overhead* demostrado por Kubernetes frente al despliegue desnudo.

\section{Trabajo Futuro}
Sugerencias para ampliar el trabajo: 
\begin{itemize}
    \item Probar aceleradores hardware en los Workers (GPU/NPU).
    \item Usar un framework de IA específico para Edge como TensorFlow Lite o ONNX Runtime.
    \item Implementar mecanismos de tolerancia a fallos automáticos en el lado de Systemd.
\end{itemize}

\vspace{1cm}

\begin{thebibliography}{99}
\bibitem{podman_docs} Red Hat, \textit{Podman Documentation}, [Online]. Disponible en: \url{https://podman.io/docs/}
\bibitem{systemd_quadlets} Freedesktop.org, \textit{systemd.service — Service unit configuration}, [Online].
\bibitem{grpc_protocol} Cloud Native Computing Foundation (CNCF), \textit{gRPC: A high performance, open source universal RPC framework}, [Online].
\bibitem{zeromq_guide} P. Hintjens, \textit{ZeroMQ: Messaging for Many Applications}, O'Reilly Media, 2013.
\bibitem{k3s_arch} Rancher Labs (SUSE), \textit{K3s Lightweight Kubernetes}, [Online]. Disponible en: \url{https://k3s.io/}
\bibitem{pytorch_edge} Paszke, A. et al., \textit{PyTorch: An Imperative Style, High-Performance Deep Learning Library}, NeurIPS 2019.
\end{thebibliography}

\end{document}